\documentclass{article}

\usepackage{booktabs}
\usepackage{multirow}
\usepackage{amsmath}
\usepackage{hyperref}
\usepackage{overpic}
\usepackage{amssymb}
\usepackage{graphicx}
\usepackage{subcaption}
\usepackage{tabularx}
\usepackage{xspace}

\DeclareMathOperator*{\argmin}{arg\,min}
\DeclareMathOperator*{\argmax}{arg\,max}
%\captionsetup[subfigure]{labelformat=simple}
\renewcommand{\thesubfigure}{\Alph{subfigure}}

\usepackage[accepted]{dlai2025}

%%% STUDENTS: FILL IN WITH YOUR OWN INFORMATION
\dlaititlerunning{ResiDual for Audio: Spectral Reweighting of Residual Streams in CLAP}

\begin{document}

\twocolumn[
%%% STUDENTS: FILL IN WITH YOUR OWN INFORMATION
\dlaititle{ResiDual for Audio: Spectral Reweighting of Residual Streams in CLAP}

\begin{center}\today\end{center}

\begin{dlaiauthorlist}
%%% STUDENTS: FILL IN WITH YOUR OWN INFORMATION
\dlaiauthor{Lorenzo Arcioni}{}
\end{dlaiauthorlist}

%%% STUDENTS: FILL IN WITH YOUR OWN INFORMATION
\dlaicorrespondingauthor{Lorenzo Arcioni}{arcioni.1885377@studenti.uniroma1.it}

\vskip 0.3in
]

\printAffiliationsAndNotice{}

\begin{abstract}
Contrastive audio-language models such as CLAP achieve remarkable zero-shot 
classification performance, yet their internal representations remain largely 
unexplored. I investigate the structure of attention head outputs in HTS-AT, 
revealing a progressive emergence of semantically specialized representations 
across the network hierarchy, tightly linked to the intrinsic dimensionality 
of per-head activations. These insights motivate the ResiDual$^*$ method, a parameter-efficient 
method that reweights the principal components of each attention head without 
modifying any pretrained weight. Evaluated under 5-fold cross-validation on 
four benchmarks, ResiDual$^*$ substantially improves over the frozen CLAP baseline.
\end{abstract}
% ------------------------------------------------------------------------------
\input{sections/Introduction}
\input{sections/Related Works}
\input{sections/Method}
\input{sections/Results}
%\input{sections/Conclusions}
%
\bibliography{references.bib}
\bibliographystyle{dlai2025}

\input{sections/Appendix}

\end{document}

